% The shield fits the header. The only decision is the jumper block.
\begin{tikzpicture}[font=\sffamily\scriptsize]

% ---- the board with the shield on it
\node[board, minimum width=54mm, minimum height=20mm] (nuc) at (0,0) {};
\node[text=white, font=\sffamily\bfseries\small] at (0,-0.55) {NUCLEO-H7A3ZI-Q};
\node[text=white, font=\tiny] at (0,0.55) {Arduino header};
\node[pinrow, minimum width=40mm, minimum height=2mm] at (0,0.95) {};

\node[hat, minimum width=46mm, minimum height=26mm] (sh) at (0,2.2) {};
\node[text=white, font=\sffamily\bfseries\small] at (0,3.15) {X-NUCLEO-IKS4A1};

% sensors on the shield
\foreach \x/\n in {-1.75/{imu\\fusion}, -0.65/{imu\\ispu}, 0.35/{accel\\mag}, 1.45/{press\\humid}}{
  \node[draw=black!70, fill=black!80, text=white, font=\tiny, align=center,
        minimum width=8mm, minimum height=5mm] at (\x,1.55) {\n};}
\node[draw=black!70, fill=gray!70, minimum width=5mm, minimum height=3mm] (elec) at (1.75,2.55) {};
\node[font=\tiny, text=white, left=0.6mm of elec] {electrode add-on};

% ---- the jumper block, which is the whole decision
\node[ext, text width=46mm, minimum height=26mm, anchor=west] (jmp) at (3.1,2.2)
  {\textbf{jumper block: three arrangements}\\[2pt]
   {\scriptsize 1. every sensor directly on the bus}\\
   {\scriptsize 2. first inertial unit as the hub}\\
   {\scriptsize 3. second inertial unit as the hub}\\[2pt]
   {\tiny positions per UM3239, read off the silkscreen}};
\draw[bus] (sh.east) -- (jmp.west);

% ---- the interrupt line, from the inertial unit down to the header
\draw[irq] (-1.75,1.3) -- node[font=\tiny, fill=white, inner sep=1pt, midway, right=0.5mm]{INT1} (-1.75,0.95);

% ---- host
\node[ext, below=12mm of nuc, text width=30mm, minimum height=9mm] (pc)
  {Host PC\\{\tiny transcript reader}};
\draw[usbcable] (nuc.south) -- (pc.north);

\node[umlnote, text width=52mm, anchor=north west] at (-2.9,-2.0)
  {One shield at a time and never two. The three shields in this inventory fit
   the same header and collide on bus addresses and on the 3.3 V budget. Two
   devices answering at one address give readings that are plausible, stable and
   wrong.};
\node[umlnote, text width=46mm, anchor=north west] at (3.1,-2.0)
  {The hub arrangements reduce the number of transactions the processor performs
   and change which device's interrupt line is the one that matters.};
\end{tikzpicture}
